\documentclass[11pt]{article}

% 한글 및 수식/이미지/레이아웃 관련 패키지
\usepackage{kotex}
\usepackage{amsmath, amssymb}
\usepackage{graphicx}
\usepackage{subcaption}
\usepackage{titlesec}
\usepackage{geometry}
\usepackage{setspace}

% 문서 설정
\geometry{a4paper, left=20mm, right=20mm, top=25mm, bottom=25mm}
\setstretch{1.31}

% 장/절 제목 서식
\titleformat{\section}
  {\centering\normalfont\large\sffamily\bfseries}
  {\thesection.}{1em}{}
\titleformat{\subsection}
  {\normalfont\normalsize\sffamily\bfseries}
  {\thesubsection.}{1em}{\itshape}

\renewcommand{\thesection}{\Roman{section}}
\renewcommand{\thesubsection}{\arabic{subsection}}

\begin{document}

%%%%%%%%%%%%%%%%%% 표지 %%%%%%%%%%%%%%%%%%
\thispagestyle{empty}
\vspace*{2cm}

\begin{center}
  {\LARGE \bf 서울과학고등학교 졸업논문}\\[30pt]
  {\Huge \bf 자기장 센서를 활용한 자기부상\\[12pt] 열차 모델에서의 저마찰 구간 최적화}\\[4cm]
  {\huge \bf 2026년 7월}\\[2cm]
  {\LARGE \bf 과학영재학교}\\[12pt]
  {\huge \bf 서울과학고등학교}\\[1cm]
  {\huge \bf 2~4~0~0~0}\\[12pt]
  {\huge \bf 김~~영~~재}
\end{center}

\newpage

%%%%%%%%%%%%%%%%%% 요약 %%%%%%%%%%%%%%%%%%
\setcounter{page}{1}

\begin{center}
  {\LARGE \bf 요 약}\\[12pt]
\end{center}

\begin{flushleft}
초록 내용을 작성한다.
\end{flushleft}

\vspace{2.5cm}

\begin{flushleft}
\textbf{주요어} : 자기부상열차, 자기장 센서, 저마찰, 최적화, 부상 동역학\\
\textbf{학번} : 24000
\end{flushleft}

\newpage

%%%%%%%%%%%%%%%%%% 본문 시작 (서론 예시) %%%%%%%%%%%%%%%%%%
\begin{center}
  {\LARGE \bf 제 1 장 서 론}\\[12pt]
\end{center}

\begin{flushleft}  
{\Large \bf 제 1 절 연구의 필요성}\\[12pt]
\end{flushleft} 

\begin{flushleft}
본문 내용을 작성한다.
\end{flushleft}

\newpage

\begin{center}
  {\LARGE \bf 제 2 장}\\[12pt]
\end{center}

\begin{flushleft}  
{\Large \bf 제 1 절}\\[12pt]
\end{flushleft} 

\begin{flushleft}
본문 내용을 작성한다.
\end{flushleft}

\newpage

%%%%%%%%%%%%%%%%%% 참고문헌 %%%%%%%%%%%%%%%%%%
\section*{참고문헌}
\addcontentsline{toc}{section}{참고문헌}

\begin{flushleft}
강서울, 정종로, 최혜화 (2024). 초고속 자기부상 열차용 자기 플럭스 제어 장치. 한국전동차학회지, 1(1), 101-110. \\
Kang, Seoul., Jeong, J. R., \& Choi, H. H. (2024). Magnetic Flux Control Device for High-Speed Maglev Trains. Journal of Korean Electric Train Society, 1(1), 101-110.
\end{flushleft}

\newpage

%%%%%%%%%%%%%%%%%% Abstract (영문 초록) %%%%%%%%%%%%%%%%%%
\thispagestyle{empty}
\vspace*{0.1cm}

\begin{center}
{\huge \bf Abstract}\\[12pt]
{\huge \bf Optimization of Low-Friction Sections in a Maglev Train Model Using Magnetic Field Sensors}\\[12pt]
\end{center}

\vspace{1.0cm}

\begin{flushright}
{\Large \bf Young Jae Kim}\\[12pt]
{\Large \bf Seoul Science High School}
\end{flushright}

\vspace{1.5cm}

\begin{flushleft}
Write an abstract in English.

\vspace{1.2cm}

\textbf{Keywords} : Maglev train, Magnetic field sensors, Low-friction, Optimization, Levitation dynamics\\
\textbf{Student Number} : \textbf{24000}
\end{flushleft}

%%%%%%%%%%%%%%%%%%
\end{document}
